\chapter{Conclusiones y trabajo futuro}
\label{cap:conclusiones}

\section{Conclusiones}

Esta tesis partió de una premisa que la literatura respalda: predecir mortalidad
por COVID-19 con datos abiertos mexicanos es un problema resuelto. Sobre esa base
desplazó la pregunta de si el modelo acierta a si acierta por igual, tomando la
entidad federativa como atributo sensible. Las conclusiones son cinco.

\textbf{Primera: la inequidad geográfica existe y es sustantiva.} Un modelo con
AUROC de \aurocGlobal{} reparte su sensibilidad de forma muy desigual entre
estados: la brecha es \gapTPRentidad{} en un punto de operación calibrado para
detectar el \sensObjetivo\,\% de las defunciones a escala nacional. Traducido, el
modelo detecta cerca de \entPeorPct{} de cada 100 defunciones en la entidad peor
servida frente a \entMejorPct{} en la mejor. La disparidad es estadísticamente
significativa ($p=\pTpr$ frente a la nula por permutación), aunque conviene
descontar lo que el azar aporta con 32 grupos desiguales: el exceso atribuible a
la entidad es de unos \excesoTpr. Y ocurre \emph{sin} que la entidad sea un
predictor del modelo, lo que descarta la solución ingenua de «no usar la
geografía».

\textbf{Segunda: no se explica por composición demográfica.} Al restringir el
análisis a adultos de 60 años o más, la brecha no solo persiste sino que se
amplía. La inequidad es una propiedad del modelo operando sobre poblaciones
heterogéneas, no un artefacto de la estructura de edades de cada estado.

\textbf{Tercera: en la población completa la disparidad reside en el umbral, no
en el ordenamiento, y eso determina qué mitigación funciona.} El AUROC por entidad
es alto en todas y la calibración por entidad es buena (error máximo
\errCalEntidad); lo que falla es aplicar un corte único a distribuciones de riesgo
distintas. La salvedad importa: entre los pacientes hospitalizados --la población
donde el desenlace es posible-- la relación se invierte, con la brecha de AUROC en
\hospGapAUROC{} y la de sensibilidad en \hospGapTPR, y en el subgrupo de 60 años o
más la dispersión del AUROC llega a \sesentaRango. El diagnóstico, y con él la
elección de mitigación, debe hacerse sobre la población en la que se desplegará. Por eso la
reponderación previa al entrenamiento no cambia nada, la restricción de
\emph{equalized odds} durante el entrenamiento resulta inestable con 32 grupos
desiguales, y solo el ajuste de umbrales por entidad reduce la brecha --de
\gapBase{} a \gapPost, un \reduccionPost\,\%-- y lo hace \emph{sin} costo de
desempeño agregado (exactitud balanceada de \accBase{} a \accPost). El costo está
en otra parte y es el que predice la teoría: al igualar la sensibilidad, la brecha
de falsos positivos crece de \postGapFPRbase{} a \postGapFPR. La mitigación no
elimina la inequidad; elige cuál de las dos tasas se reparte por igual. La
recomendación práctica que se desprende es diagnosticar primero dónde vive el
sesgo --brechas de AUROC frente a brechas de sensibilidad-- y elegir la familia de
mitigación en consecuencia.

\textbf{Cuarta: la inequidad no es un problema de cobertura de datos.} Al excluir
por completo una entidad del entrenamiento y evaluar sobre ella, el desempeño
apenas se mueve --pérdida mediana de AUROC \loeoMediana{} sobre \loeoNent{}
entidades-- y las brechas no se acentúan. La pérdida además no guarda relación con
el tamaño de la entidad ($r=\loeoCorrTamr$, $p=\loeoCorrTamp$). Junto con el
fracaso de la reponderación, esto descarta la explicación más intuitiva: el modelo
no falla en unos estados por haber visto pocos datos suyos; funciona igual sin
haber visto ninguno. Tampoco es un efecto de la práctica de registro de la
variable dominante: al reentrenar sin \texttt{\impPrimeraVar} la brecha de AUROC
no se reduce, sino que sube a \sinNeuGapAUROC. Lo que produce la inequidad es
aplicar una regla común a poblaciones distintas.

\textbf{Quinta: la separación metodológica no es un detalle.} Dos decisiones de
diseño cambian conclusiones sustantivas. Auditar la calibración por grupo sobre
puntajes sin calibrar habría reportado un error de \errCalCrudo{} y una conclusión de
«calibración inequitativa» que corresponde en realidad al sesgo global de
entrenamiento; medido sobre riesgo calibrado, el error es \errCalEntidad{} y la
calibración por entidad es buena. Y estimar los umbrales de mitigación sobre el
propio conjunto de prueba habría llevado a declarar la brecha prácticamente
eliminada (\gapPostOpt) cuando la cifra desplegable es \gapPost, \factorOptimismo{} veces
mayor.

Estas conclusiones se sostienen dentro de los límites declarados en la
sección~\ref{sec:alcance} y con las limitaciones de la
sección~\ref{sec:limitaciones}. En particular, la hipótesis~\ref{h2} queda
refutada en la forma en que se operacionalizó --el índice de marginación estatal
no predice el sesgo-- y la hipótesis~\ref{h4} sin poner a prueba, si bien la
generalización geográfica interna, que es lo más cercano que puede evaluarse sin
una segunda fuente, resultó buena.

\section{Trabajo futuro}
\label{sec:trabajo_futuro}

Ordenado por prioridad, entendida como la capacidad de cada tarea para modificar
las conclusiones anteriores:

\begin{enumerate}[leftmargin=*]
  \item \textbf{Extender la cuantificación de incertidumbre.} El eje geográfico
    ya cuenta con intervalos de confianza y prueba de permutación
    (sección~\ref{sec:incertidumbre_res}); falta hacer lo propio con los subgrupos
    interseccionales, donde los tamaños son mucho menores y la precisión es la
    principal duda abierta.
  \item \textbf{Indicadores de infraestructura de salud.} El índice de
    marginación de CONAPO ya está incorporado y no explica el sesgo. El paso
    siguiente es sustituirlo por el mecanismo que \ref{h2} postula en realidad:
    camas, personal y equipamiento por habitante, idealmente por unidad médica y
    no por entidad.
  \item \textbf{Validación externa en la base DGIS.} Es el paso que pone a prueba
    \ref{h4} y el que la literatura casi nunca da. Tiene un prerrequisito que este
    trabajo dejó identificado: SAEH es una base de egresos y su población son
    pacientes hospitalizados, con una prevalencia del desenlace cercana a
    \hospPrev{} frente a \cohorteTasaDef\,\% aquí. Debe validarse el modelo
    restringido a hospitalizados (sección~\ref{sec:sensibilidad_res}); usar el de la
    cohorte completa haría pasar un cambio de espectro por un fallo de
    transportabilidad. Requiere además reconstruir las once comorbilidades desde
    los códigos CIE-10 de SAEH, lo que añade desplazamiento de medición al de
    población.
  \item \textbf{Validación temporal.} Dividir por olas epidémicas y entrenar en un
    año para evaluar en otro, con los cortes anuales de la DGE. La generalización
    geográfica ya se evaluó (sección~\ref{sec:loeo_res}); la temporal no.
  \item \textbf{Umbrales por grupo con regularización.} Contraer los umbrales de
    entidades pequeñas hacia el umbral global, en lugar de estimarlos de forma
    independiente. Es la vía natural para atacar la brecha residual de \gapPost{}
    sin perseguir ruido.
  \item \textbf{Sesgo de verificación.} Modelar explícitamente la propensión a la
    detección por entidad, para separar la inequidad del modelo de la inequidad del
    proceso de captura.
  \item \textbf{Desenlace de hospitalización y utilidad clínica.} Extender el
    análisis al segundo desenlace y añadir análisis de curva de decisión
    (\emph{net benefit})~\parencite{vickers2006} y explicaciones locales, previstos en
    el protocolo y no ejecutados.
\end{enumerate}
